% !TEX TS-program = pdflatexmk
\documentclass[12pt]{amsart}

%\usepackage[parfill]{parskip}    % Activate to begin paragraphs with an empty line rather than an indent

\usepackage[margin=1in]{geometry}

\usepackage{amsmath,amssymb,amsthm,latexsym,graphicx}
\usepackage[normalem]{ulem}
\usepackage{setspace} %used for doublespacing, etc.
\usepackage{hyperref}
\usepackage[shortlabels]{enumitem}
\usepackage{cancel}
\usepackage[dvipsnames,usenames]{color}
\usepackage[all]{xy}
\DeclareGraphicsRule{.tif}{png}{.png}{`convert #1 `dirname #1`/`basename #1 .tif`.png}

%Some useful environments.
\newtheorem{theorem}{Theorem}
\newtheorem{corollary}[theorem]{Corollary}
\newtheorem{conjecture}[theorem]{Conjecture}
\newtheorem{lemma}[theorem]{Lemma}
\newtheorem{proposition}[theorem]{Proposition}
\newtheorem{definition}[theorem]{Definition}
\newtheorem{example}[theorem]{Example}
\newtheorem{axiom}{Axiom}
\theoremstyle{remark}
\newtheorem{remark}{Remark}
\newtheorem*{exercise}{Exercise}%[section]

%Some useful shortcuts for our favorite fields
\newcommand{\RR}{{\mathbb R}}
\newcommand{\NN}{{\mathbb N}}
\newcommand{\ZZ}{{\mathbb Z}}
\newcommand{\QQ}{{\mathbb Q}}
\newcommand{\CC}{{\mathbb C}}

%Some useful shortcuts for formatting lists
\newcommand{\bc}{\begin{center}}
\newcommand{\ec}{\end{center}}
\newcommand{\be}{\begin{enumerate}}
\newcommand{\ee}{\end{enumerate}}
\newcommand{\bi}{\begin{itemize}}
\newcommand{\ei}{\end{itemize}}

%Some useful shortcuts for formatting mathematical symbols
\newcommand{\ol}[1]{\overline{#1}}
\newcommand{\oimp}[1]{\overset{#1}{\Longleftrightarrow}} %labeled iff symbol
\newcommand{\bv}[1]{\ensuremath{ \vec{\mathbf{#1}}} } %makes a vector.
\newcommand{\mc}[1]{\ensuremath{\mathcal{#1}}} %put something in caligraphic font
\newcommand{\mpg}[1]{\marginpar{ #1}} %to put comments in margins
\newcommand{\bsl}[1]{\texttt{\symbol{92}{\em #1}}} %for backslashes.
\newcommand{\normale}{\trianglelefteq}
\newcommand{\normal}{\triangleleft}

%Commenting tools
\usepackage{soul}
\definecolor{highlight}{rgb}{1,0.6,0.6}
\sethlcolor{highlight}
\newcommand{\hlm}[1]{\colorbox{highlight}{$\displaystyle #1$}}
\newtheoremstyle{mycomment}{\topsep}{-0in}{\small \itshape \sffamily}{}{\small \itshape\sffamily}{:}{.5em}{}
\theoremstyle{mycomment}
\newtheorem*{acomment}{\color{BrickRed}{Comment}}
\newcommand{\com}[1]{{\color{OliveGreen}\begin{acomment}{#1} %#2 \color{black} 
\end{acomment}\noindent}}
%\newcommand{\com}[1]{{\color{BrickRed}{\\ Comment:}\color{OliveGreen}{#1} \\}}
\newcommand{\red}[1]{{\color{BrickRed} #1}}
\def\midd{\,\middle\vert\, }
\newcommand{\blue}[1]{{\color{MidnightBlue}#1}}
\newcommand{\green}[1]{{\color{OliveGreen}#1}}
\newcommand{\mwrong}[2]{\red{\cancel{#1}}\green{#2}}
\newcommand{\wrong}[2]{\red{\sout{#1}}\green{#2}}
\definecolor{OliveGreen}{rgb}{.3,.5,.2}
\definecolor{MidnightBlue}{rgb}{.3,.4,.6}

\title{Assignment \#10 \today}

\author{Coordinator: Glen Woodworth}
\begin{document}
\maketitle
%\setstretch{2.5} %Use for 2.5 spacing
\doublespacing %Use for double spacing


\section*{Section 8.1}

\begin{exercise}[\textbf{8.1.1} --- Stefano Fochesatto] For each of the following pairs of integers $a$ and $b$, determine their greatest common divisor $d$
  and write $d$ as a linear combination $ax + by$ of $a$ and $b$. 
  \begin{enumerate}
    \item[a.] $a = 20$, $b = 13$ 
    \begin{proof} Using the euclidean algorithm we get the following, 
      \begin{align*}
        20 &= 13(1) + 7\\
        13 &= 7(1) + 6\\
        7 &= 6(1) + 1\\
        6 &= 6(1) + 0
      \end{align*}
      Which gives us $(20, 13) = 1$. By back substitution we get the following, $1 = 7 - 6 = 7 - (13 - 7) = 2(7) - 13 = 2(20 - 13) - 13 = 2(20) - 3(13)$.
    \end{proof}
    \vspace{.15in}


    \item[c.] $a = 11391$, $b = 5673$
    \begin{proof} Using the euclidean algorithm we get the following, 
      \begin{align*}
        11391 &= 5673(2) + 45\\
      5673 &= 45(126) + 3\\
        45 &= 3(15) + 0
      \end{align*}
      Which gives us $(11391, 5673) = 3$. By back substitution we get the following, $3 = 5673 - 45(126) = 5673 - (11391 - (2)5673)(126) = 253(5673) - 126(11391)$
    \end{proof}
    \vspace{.15in}



    \item[e.] $a = 91442056588823$, $b = 779086434385541$
    \begin{proof} Using the euclidean algorithm we get the following, 
      \begin{align*}
        779086434385541 &= 91442056588823(8) + 47549981674957\\
        91442056588823 &= 47549981674957(1) + 43892074913866\\
        47549981674957 &= 43892074913866(1) + 3657906761091\\
        43892074913866 &= 3657906761091(11) + 3655100541865\\
        3657906761091 &=  3655100541865(1) + 2806219226\\
        3655100541865 &= 2806219226(1302) + 1403109613\\
        2806219226 &= 1403109613(2) + 0
      \end{align*}
      Which gives us $(91442056588823,779086434385541) = 1403109613$. By back substitution we get, $1403109613 = 277522(91442056588823) - 32573(779086434385541)$
    \end{proof}

  \end{enumerate}
\end{exercise}
\vspace{.5in}



\begin{exercise}[\textbf{8.1.7} --- Stefano Fochesatto] Find the generator for the ideal $(85, 1 + 13i)$ in $\ZZ[i]$, i.e. the greatest common divisor for $85$
  and $1 + 13i$, by the Euclidean Algorithm. Do the same for the ideal $(47 - 13i, 53 + 56i)$.
    \begin{proof}[Solution] As outlined in Example $3$ of Section 8.1 we know that The Gaussian integers are a Euclidean Domain under the norm $N(a + bi) = a^2 + b^2$. 
      Recall the division algorithm for The Gaussian integers under norm $N$, for all $\alpha, \beta \in \ZZ[i]$ there exists some $(p + qi), \gamma \in \ZZ[i]$ such that 
      $\alpha = (p + qi)\beta + \gamma$ where $N(\gamma) \leq \frac{1}{2}N(\beta)$. 
      We begin the Euclidean algorithm, note that 
      \begin{equation*}
        \frac{85}{1 + 13i} = \frac{85(1 - 13i)}{(1 + 13i)(1 - 13i)} = \frac{85 - (85)13i}{170} = \frac{1}{2} - \frac{13}{2}i. 
      \end{equation*}
      Choosing a close $(p + qi) \in \ZZ[i]$ we get $(p + qi) = 1 - 7i$. Finding the product $(1 + 13i)(1 - 7i) = 92 + 6i$. Solving for $\gamma = 85 - (92 + 6i) = -7 - 6i$.
      Double checking norms we get, $N(-7 - 6i) = 7^2 + 6^2 = 85 \leq 85 = \frac{1}{2}170 N(1 + 13i)$. 
      So our first step looks like, 
      \begin{equation*}
        85 = (1 - 7i)(1 + 13i) + (-7 - 6i).
      \end{equation*}
      Computing the next step we first check the quotient,
      \begin{equation*}
        \frac{1 + 13i}{-7 - 6i} = \frac{(1 + 13i)(-7 + 6i)}{(-7 - 6i)(-7 + 6i)} = \frac{-85 - 85i}{85} = -1 - 1i.
      \end{equation*}
      Which gives us our final step, 
      \begin{equation*}
        1 + 13i = (-7 - 6i)(-1 - 1i) + 0.
      \end{equation*}
      Therefore the greatest common divisor for $85$ and $1 + 13i$ is $(-7 - 6i)$.\\\\

      
      Applying the Euclidean algorithm for $47 - 13i$ and $53 + 56i$, first we compute the following quotient, 
      \begin{equation*}
        \frac{53 + 56i}{47 - 13i} = \frac{(53 + 56i)(47 + 13i)}{(47 - 13i)(47 + 13i)} = \frac{1763 + 3321i}{2378} = \frac{43}{58} + \frac{81}{58}i 
      \end{equation*}
      Choosing a close $(p + qi) \in \ZZ[i]$ we get $(p + qi) = 1 + i$. Finding the product $(47 - 13i)(1 + i) = 60 + 34i$. Solving for $\gamma = (53 + 56i) - (60 + 34i) = -7 + 22i$.
      Double checking norm we get, $N(-7 + 22i) = 533 \leq 1189  = \frac{2378}{2} = \frac{1}{2}N(47 - 13i)$. So our first step looks like, 
      \begin{equation*}
        53 + 56i = (1 + i)(47 - 13i) + (-7 + 22i).
      \end{equation*}
      Computing the next step we first check the quotient, 
      \begin{equation*}
        \frac{47 - 13i}{-7 + 22i} = \frac{(47 - 13i)(-7 - 22i)}{(-7 + 22i)(-7 - 22i)} = \frac{-615 - 943i}{533} = -\frac{15}{13} - \frac{23}{13}i.
      \end{equation*}
      Choosing a close $(p + qi) \in \ZZ[i]$ we get $(p + qi) = -1 - 2i$. Finding the product $(-7 + 22i)( -1 - 2i) = 51 - 8i$. Computing $\gamma = (47 - 13i) - (51 - 8i) = - 4 -5i$.
      So our next step is 
      \begin{equation*}
        47 - 13i = (-1 - 2i)(-7 + 22i) + (-4-5i).
      \end{equation*}
      For the next step we compute the following quotient, 
      \begin{equation*}
        \frac{-7 + 22i}{-4-5i} = \frac{(-7 + 22i)(-4+5i)}{(-4-5i)(-4+5i)} = \frac{-82 - 123 i}{41} = -2 -3i.
      \end{equation*}
      Which gives us our final step, 
      \begin{equation*}
        -7 + 22i = (-2 -3i)(-4-5i) + 0.
      \end{equation*}
      Therefore greatest common divisor for $47 - 13i$ and $53 + 56i$ is $(-4-5i)$.

    \end{proof}

\end{exercise}
\vspace{.5in}

\section*{Section 8.2}


\begin{exercise}[\textbf{8.2.1} --- Stefano Fochesatto]  Prove that in an Principal Ideal Domain $R$, two ideals $(a)$ and $(b)$ 
  are comaximal if and only if a greatest common divisor of $a$ and $b$ is $1$. 
    \begin{proof} Suppose $R$ is a Principal Ideal Domain and two ideal $(a)$ and $(b)$ 
      are comaximal. By the definition of comaximal we know that $(a) + (b) = R$. By Proposition 6 we know that if 
      $d$ is a generator for the principal ideal generated by $a$ and $b$ then $d = \gcd(a, b)$. Therefore $(a), (b) \subseteq (d)$ and 
      by (ii) of Proposition 6, $(a) + (b) \subseteq (d)$. Since $(a)$ and $(b)$ are comaximal we know that $(a) + (b) = R$ therefore $(d) = R$. 
      Thus $d = 1$ and therefore $gcd(a, b) = 1$. 

      Suppose $R$ is a Principal Ideal Domain with elements $a$ and $b$ such that $\gcd(a, b) = 1$.  Consider 
      $(a)$ and $(b)$ by (ii) of Proposition 6 we know that $(1) = (a)x + (b)y$ where $x, y \in R$. Thus $(a)$ and $(b)$
      are comaximal. 
    

    \end{proof}

\end{exercise}
\vspace{.5in}





\end{document}